\section{The typed dependency graph}
\label{sec:dependency-graph}

The principal completion chain is
\begin{equation}
\begin{gathered}
\text{TPC-73: primitive support}\\[-0.1em]
\downarrow\\[-0.1em]
\text{TPC-74: direct two-corner entries}\\[-0.1em]
\downarrow\\[-0.1em]
\text{TPC-75: completion duality}
\longrightarrow
\text{TPC-76: provenance lifts}\\[-0.1em]
\longrightarrow
\text{TPC-77: sparse components}
\longrightarrow
\text{TPC-78: complex-scalar forest and motif reductions}.
\end{gathered}
\label{eq:completion-chain}
\end{equation}
Two auxiliary branches leave this chain.  A verified anchor synthesis
and a direct synthesis can enter TPC-79 only if they share one target:
\begin{equation}
\left.
\begin{array}{c}
\text{direct synthesis into \(V\)}\\
\text{anchor synthesis into the same \(V\)}
\end{array}
\right\}
\longrightarrow
\text{TPC-79: hybrid/orphan quotient}.
\label{eq:anchor-branch}
\end{equation}
Likewise, the TPC-80 theorem requires a literal scalar Hermitian
matrix \(H_X^{\rm hol}\) and a nonnegative majorant
\(P_X^{\rm hol}\) on the same row graph:
\begin{equation}
(H_X^{\rm hol},P_X^{\rm hol})
\text{ on one scalar actual row graph}
\longrightarrow
\text{TPC-80: Perron/holonomy audit}.
\label{eq:holonomy-branch}
\end{equation}

None of \eqref{eq:completion-chain}--\eqref{eq:holonomy-branch} is a
determinant-frequency statement.  The only typed route from those
completion branches to that statistic in the present block is
\begin{equation}
\begin{gathered}
\text{literal pre-bin coefficient \(\cB_Xz_X\)}
\xrightarrow{\ \cJ_X\ }
a_X=(A_{C_X}(n))_n\\
\downarrow\\[-0.2em]
(Z_X,D_X)
\longrightarrow
\mathfrak F_{h_0}(X)
\longrightarrow
\text{the conditional TPC-32 threshold}.
\end{gathered}
\label{eq:determinant-chain}
\end{equation}
Here \(\cJ_X\) is the exact determinant dictionary, not a label
identification.

\begin{theorem}[Typed dependency theorem]
\label{thm:typed-dependency}
The implications in \eqref{eq:completion-chain} are composable only
through their stated synthesis and dual maps.  The TPC-79 branch
additionally requires a common target and separable native variables;
the TPC-80 branch additionally requires scalar Hermitian entries and
one common actual row graph.  A conclusion from either branch enters
\eqref{eq:determinant-chain} only through an explicit estimate for the
literal \(\cB_X\), the verified binning map \(\cJ_X\), and its
dictionary residual.  Shared labels, shared support, or a numerical
correlation do not replace these maps.
\end{theorem}

\begin{proof}
The completion results of TPC-73--TPC-78 concern representations and
norms of synthesis targets
\citep{WangTPC73,WangTPC74,WangTPC75,WangTPC76,WangTPC77,WangTPC78}.
TPC-79 forms an infimal convolution in one target space
\citep{WangTPC79}.  TPC-80 compares quadratic forms of two scalar
matrices on one vertex space \citep{WangTPC80}.  By contrast,
\(a_X\) is indexed by determinant bins.  These domains and codomains
are not canonically identical.  Composition is therefore defined
only after the intervening maps and their literal multipliers are
specified.  TPC-81 makes precisely that extra map explicit
\citep{WangTPC81}.  Without it, the displayed objects are merely
different typed data sets.
\end{proof}

\begin{remark}[No silent Perron-pair identification]
The TPC-80 pair
\((H_X^{\rm hol},P_X^{\rm hol})\) and the TPC-81 normalized
missing-Gram pair \((R_{\cM,X},P_{\cM,X})\) have different
provenance.  TPC-80 can sharpen the missing-Gram estimate only after
an exact equality or a norm-preserving intertwining between these
pairs is proved on the same actual scalar graph.
\end{remark}

\begin{remark}[Direction of inference]
The graph records legal implication, not automatic quantitative
success.  An exact TPC-78 reduction may return a large optimum; an
exact TPC-79 dual witness may certify no saving in the declared
hybrid class; and a TPC-80 equality certificate may show exact Perron
saturation.  These are useful stop certificates but they do not close
\eqref{eq:determinant-chain}.
\end{remark}
